% ============================================================================
% Template da Fase 3 — Desenvolvimento dos Resultados Práticos
% Disciplina: MO810A/MC959A — Sistemas Multiagentes com Modelos de Linguagem
% Instituto de Computação — UNICAMP (1º Semestre de 2026)
% ============================================================================
\documentclass[12pt,a4paper]{article}

\usepackage[utf8]{inputenc}
\usepackage[T1]{fontenc}
\usepackage[brazil]{babel}
\usepackage{geometry}
\geometry{margin=2.5cm}

\usepackage{setspace}
\onehalfspacing

\usepackage{titlesec}
\usepackage{hyperref}
\usepackage{graphicx}
\usepackage{enumitem}
\usepackage{xcolor}
\usepackage{fancyhdr}
\usepackage{amsmath,amssymb}
\usepackage{booktabs}
\usepackage{array}
\usepackage{longtable}
\usepackage[numbers]{natbib}
\newcommand{\citeonline}[1]{\citeauthor{#1}~\cite{#1}}

% --- Cabeçalho e Rodapé ---
\setlength{\headheight}{14.5pt}
\pagestyle{fancy}
\fancyhf{}
\rhead{\footnotesize MO810A/MC959A -- Sistemas Multiagentes com LLMs}
\lhead{\footnotesize Fase 3 -- Resultados Práticos}
\cfoot{\thepage}

% --- Estilo para caixas de instrução ---
\usepackage{mdframed}
\newmdenv[
  linecolor=gray!60,
  backgroundcolor=gray!8,
  roundcorner=4pt,
  innertopmargin=8pt,
  innerbottommargin=8pt,
  innerleftmargin=10pt,
  innerrightmargin=10pt,
  skipabove=6pt,
  skipbelow=6pt
]{instrucao}

\newcommand{\dica}[1]{%
  \begin{instrucao}
  \small\textcolor{gray!80!black}{\textbf{Instruções:} #1}
  \end{instrucao}
}

\begin{document}

\begin{center}
  {\LARGE \textbf{Universidade Estadual de Campinas}}\\[0.2cm]
  {\large Instituto de Computação}\\[0.5cm]
  {\large \textbf{MO810A / MC959A -- Tópicos em Inteligência Artificial:}}\\[0.1cm]
  {\large \textbf{Sistemas Multiagentes com Modelos de Linguagem}}\\[0.2cm]
  {\normalsize Prof.\ Dr.\ Julio Cesar dos Reis --- 1º Semestre de 2026}\\[0.6cm]
  \rule{\textwidth}{0.4pt}\\[0.4cm]
  {\Large \textbf{Fase 3 -- Desenvolvimento dos Resultados Práticos}}\\[0.3cm]
  {\large \textbf{[Título do Projeto]}}\\[0.3cm]
  \rule{\textwidth}{0.4pt}\\[0.4cm]
\end{center}

\noindent
\begin{tabular}{@{}ll@{}}
  \textbf{Integrante 1:} & Nome Completo --- RA XXXXXX \\
  \textbf{Integrante 2:} & Nome Completo --- RA XXXXXX \\
  \textbf{Repositório Git:} & \url{https://github.com/usuario/repositorio} \\
  \textbf{Data de Entrega:} & \today \\
\end{tabular}

\vspace{0.6cm}

% ============================================================================
\section{Visão Geral da Implementação e Arquitetura em Operação}
\label{sec:visao-geral}
\dica{Apresente o estado operacional do Sistema Multiagente com Modelos de Linguagem (LLMs) especificado na Fase 2, destacando como os agentes, a orquestração, o uso de ferramentas (\textit{tool use}) e o acesso às fontes de dados funcionam de ponta a ponta.}

% ============================================================================
\section{Tecnologias, Frameworks e Infraestrutura de Execução}
\label{sec:tecnologias}
\dica{Descreva detalhadamente as tecnologias e os \textit{frameworks} utilizados: bibliotecas de construção e orquestração de agentes, modelos de linguagem empregados com provedor e versão, bancos de dados e índices vetoriais, APIs e serviços expostos aos agentes como ferramentas.}

% ============================================================================
\section{Especificação Técnica dos Agentes}
\label{sec:especificacao-agentes}
\dica{Detalhe as tomadas de decisão técnicas referentes aos agentes: papéis, \textit{prompts} e habilidades (\textit{skills}), contratos de entrada e saída e critérios explícitos de término.}

% ============================================================================
\section{Padrão de Orquestração, Protocolo de Mensagens e Uso de Ferramentas}
\label{sec:orquestracao-ferramentas}
\dica{Descreva o padrão de orquestração implementado, o protocolo de troca de mensagens entre agentes e serviços, os esquemas de chamada das ferramentas (\textit{tool calling}) e os mecanismos de validação de suas saídas.}

% ============================================================================
\section{Mecanismos de Memória, Limites de Execução, Custo e Tratamento de Falhas}
\label{sec:memoria-limites-falhas}
\dica{Explicite o funcionamento dos mecanismos de memória (curto e longo prazo), os orçamentos e limites de execução (passos, chamadas, \textit{tokens}, tempo), o controle de custo computacional e o tratamento de erros, reparos e recuperação de falhas.}

% ============================================================================
\section{Ajustes Realizados em Relação à Proposta Original e Justificativas}
\label{sec:ajustes}
\dica{Explicite os ajustes realizados na implementação em relação à proposta original das Fases 1 e 2 e apresente a respectiva justificativa técnica para cada mudança.}

% ============================================================================
\section{Execução Ponta a Ponta e Resultados Práticos Obtidos}
\label{sec:resultados-praticos}
\dica{Demonstre a execução ponta a ponta do sistema multiagente gerando os artefatos práticos esperados (ingestão documental, recuperação híbrida, tutoria, avaliação por rubrica, memória temporal e planejamento CP-SAT), reservando os experimentos comparativos e ablações para a Fase 4.}

% ============================================================================
\section*{Declaração de uso de IA generativa no processo de elaboração}

\bibliographystyle{plainnat}
\bibliography{references}

\end{document}

